\section{Accountability, Oversight, and Who Answers}

\subsection{The question, stated as patients state it}

The accountability concern in the surveyed literature is not that nobody is
responsible. It is that everybody is, which amounts to the same thing
\cite{Jauniaux2025, WuSemiAI2026}. When a human surgeon errs, the chain of
liability is familiar to the person harmed and to their family. When a
semi-autonomous movement or a model-generated suggestion contributes to harm,
participants report being told that the surgeon, the hospital, the software
developer, the device manufacturer, and the sponsor are all involved, which
answers nothing and leaves them with no one to call.

This protocol answers with a rule rather than a list: exactly one accountable
party per failure class, named in advance, published before consent. The rule is
checkable, because a matrix with two accountable parties in a single row would
fail it and the matrix is published. The oversight package this is drawn from is
the author's site documentation set \cite{paisitedocpk}.

\subsection{One accountable party per failure class}

Figure 29 renders the matrix. Nine failure classes, nine single accountable
assignments, and a final column giving the person the participant should
actually call. That last column is the only black fill in the figure, because it
is the only part of the matrix that is an instruction rather than a description.

Two accountable parties in one cell would be the same as none. The invariant is
stated in the source and verified: nine rows, nine A assignments, no row with
two and no row with zero.

\begin{pafig}[d2-type: grid matrix, responsibility assignment]
\node[d2cellh,minimum height=9mm,text width=44mm,align=left] (h0) at (0,0) {If this goes wrong};
\foreach \i/\lab in {1/{Operating\\surgeon},2/{Sponsor-\\Investigator},3/{Site PI},4/{DSMB},5/{IRB and\\FDA}}{
  \node[d2cellh,minimum height=9mm,text width=17mm] at ({3.16+(\i-1)*1.83},0) {\lab};}
\node[d2cell,fill=padark,text=protowhite,font=\tiny\sffamily\bfseries,minimum height=9mm,text width=44mm]
  (h6) at (13.64,0) {You call};
% racell{column index}{row index}{letter}{style} - one matrix cell. The style
% token is supplied explicitly rather than derived, so the fill is unambiguous.
\newcommand{\racell}[4]{%
  \node[#4,minimum height=7mm,text width=17mm] at ({3.16+(#1-1)*1.83},{-0.86-(#2-1)*0.715}) {#3};}
\foreach \r/\lab/\who in {%
  1/{A motion during your operation}/{the surgeon, by name, recorded in your consent},
  2/{A device malfunction or a halt}/{the Sponsor-Investigator},
  3/{A drug dose or a restart decision}/{the Site PI first, then the Sponsor},
  4/{A serious adverse event}/{the Site PI, on the 24-hour line},
  5/{A decision to pause the study}/{you do not need to call; you are told},
  6/{Your data being mishandled}/{the Site PI, then the IRB directly},
  7/{A software version change}/{you are re-consented before anything proceeds},
  8/{A bill you did not expect}/{the study coordinator, then the Sponsor},
  9/{A complaint about how you were treated}/{the IRB directly, without going through the site}}{
  \node[d2celll,minimum height=7mm,text width=44mm] at (0,{-0.86-(\r-1)*0.715}) {\lab};
  \node[d2cell,fill=padark,text=protowhite,minimum height=7mm,text width=44mm,align=left]
    at (13.64,{-0.86-(\r-1)*0.715}) {\who};}
\tikzset{
  raA/.style={d2cell,fill=protoblue,text=protowhite,font=\tiny\sffamily\bfseries},
  raR/.style={d2cell,fill=pablue1,text=protowhite},
  raC/.style={d2cell,fill=pablue2},
  raI/.style={d2cell,fill=pagrayl}}
\racell{1}{1}{A}{raA}\racell{2}{1}{C}{raC}\racell{3}{1}{R}{raR}\racell{4}{1}{I}{raI}\racell{5}{1}{I}{raI}
\racell{1}{2}{R}{raR}\racell{2}{2}{A}{raA}\racell{3}{2}{R}{raR}\racell{4}{2}{C}{raC}\racell{5}{2}{I}{raI}
\racell{1}{3}{C}{raC}\racell{2}{3}{A}{raA}\racell{3}{3}{R}{raR}\racell{4}{3}{C}{raC}\racell{5}{3}{I}{raI}
\racell{1}{4}{R}{raR}\racell{2}{4}{A}{raA}\racell{3}{4}{R}{raR}\racell{4}{4}{C}{raC}\racell{5}{4}{I}{raI}
\racell{1}{5}{I}{raI}\racell{2}{5}{R}{raR}\racell{3}{5}{I}{raI}\racell{4}{5}{A}{raA}\racell{5}{5}{C}{raC}
\racell{1}{6}{I}{raI}\racell{2}{6}{R}{raR}\racell{3}{6}{A}{raA}\racell{4}{6}{I}{raI}\racell{5}{6}{C}{raC}
\racell{1}{7}{I}{raI}\racell{2}{7}{A}{raA}\racell{3}{7}{R}{raR}\racell{4}{7}{C}{raC}\racell{5}{7}{C}{raC}
\racell{1}{8}{I}{raI}\racell{2}{8}{A}{raA}\racell{3}{8}{R}{raR}\racell{4}{8}{I}{raI}\racell{5}{8}{I}{raI}
\racell{1}{9}{I}{raI}\racell{2}{9}{C}{raC}\racell{3}{9}{R}{raR}\racell{4}{9}{I}{raI}\racell{5}{9}{A}{raA}
\begin{scope}[on background layer]
\node[d2cont,fit={(-2.45,0.62) (16.05,-7.05)}] (M) {};
\end{scope}
\node[d2title,anchor=south west] at (M.north west) {accountability.matrix};
\legkey{-2.3}{-7.7}{protoblue}{A, accountable, exactly one per row}
\legkey{2.4}{-7.7}{pablue1}{R, responsible}
\legkey{5.6}{-7.7}{pablue2}{C, consulted}
\legkey{8.6}{-7.7}{pagrayl}{I, informed}
\legkey{11.4}{-7.7}{padark}{the call you make}
\node[font=\scriptsize\sffamily\itshape,text=protogray,anchor=north west,align=left]
  at (-2.3,-8.3) {Nine failure classes, nine single accountable assignments. Two accountable parties in one cell would be the same as none, which is\\the failure the concern describes. The far-right column is the only black fill in the figure, because it is the only instruction in it.};
\end{pafig}
\vspace{-0.7cm}
\figcaption{Figure 29. The responsibility matrix with exactly one accountable party in every\\
row and a final column, the only black fill in the figure, naming the person the\\
participant should actually call when that particular thing goes wrong for them.}

\vspace{0.30em}

\noindent\begin{tabularx}{\textwidth}{L{4.6cm} L{3.7cm} Y}
\toprule
\textbf{If this goes wrong} & \textbf{Accountable party} & \textbf{Who you call} \\
\midrule
A motion during your operation & Operating surgeon & The surgeon, by name, recorded in your consent \\
A device malfunction or a halt & Sponsor-Investigator & The Sponsor-Investigator \\
A drug dose or a restart decision & Sponsor-Investigator & The Site PI first, then the Sponsor \\
A serious adverse event & Sponsor-Investigator & The Site PI, on the 24-hour line \\
A decision to pause the study & DSMB & You do not need to call; you are told \\
Your data being mishandled & Site PI & The Site PI, then the IRB directly \\
A software version change & Sponsor-Investigator & You are re-consented before anything proceeds \\
A bill you did not expect & Sponsor-Investigator & The study coordinator, then the Sponsor \\
A complaint about how you were treated & IRB & The IRB directly, without going through the site \\
\bottomrule
\end{tabularx}

\vspace{0.30em}

\subsection{The oversight bodies, and what each may actually do}

Authority and advice are different things, and a governance section that lists
bodies without distinguishing them is not informative. The strongest action each
body may take is stated below.

\vspace{0.30em}

\noindent\begin{tabularx}{\textwidth}{L{4.1cm} L{4.2cm} Y}
\toprule
\textbf{Body} & \textbf{Independent of the sponsor} & \textbf{Strongest action it may take} \\
\midrule
Data and Safety Monitoring Board & Yes & Halt the study at a cohort boundary, or immediately \\
Physical AI Safety Review Committee & Yes & Require re-gating and a Unified Safety Level reassessment \\
Institutional Review Board & Yes & Suspend or withdraw approval; receives your complaints directly \\
Contract research organisation monitor & Contracted by the sponsor & Raise a finding to closure; cannot stop a procedure \\
Independent safety monitor, in the room & Yes, per procedure & Stop the procedure; cannot change the protocol \\
Operating surgeon & No & Stop, convert, or decline any model-proposed motion \\
Sponsor-Investigator & No & Pause enrollment; cannot overrule a DSMB halt \\
\bottomrule
\end{tabularx}

\vspace{0.30em}

\noindent The two rows a participant should note are the last two. The
Sponsor-Investigator, who is accountable for six of the nine failure classes,
cannot overrule a halt by the Data and Safety Monitoring Board. And the
independent safety monitor, who can stop any procedure, cannot change the
protocol. Neither of those constraints would be visible in a list of bodies.

\subsection{Team experience, disclosed rather than implied}

The concern reported in the literature is not general surgical experience; it is
experience with this exact platform, this software version, this autonomy level,
and this operation \cite{WuSemiAI2026, Jauniaux2025}. General experience is what
gets offered, and it is not the same thing.

Four figures are published to the participant before consent. The operating
surgeon's total robotic pancreatoduodenectomy case count. The site's position on
the learning curve defined by the nationwide comparator \cite{DutchCohort2025}.
The number of procedures the operating surgeon has performed on this specific
platform at this specific autonomy level. And whether a trainee will participate
in any part of the procedure, and if so which part.

Case count is the right question because the comparator series shows outcome
varying with it, which \S 10.2 quantifies. A participant who is told only that
their surgeon is experienced has been given an adjective; a participant who is
told a number can compare it to the phase boundaries in that table.

\subsection{The audit trail, and your right to read it}

Every model advisory, every gate decision, every surgeon approval, and every
safety-limit event is written to a hash-chained record bound to a deterministic
seed and a specific repository commit, under 21 CFR part 11
\cite{ecfr2026part11, cfr312paiuni}. The chain is append-only: an entry cannot
be altered after it is written without breaking the hash of every entry after
it, which is the property that makes requesting an extract worth doing.

A participant may request their own extract at any time. The extract contains
every advisory issued during their procedure with its confidence value, every
gate decision including every rejection and the limit that caused it, every
surgeon approval with its timestamp and signature, and every safety-limit event
with the measured value. It does not contain any other participant's data, the
model weights, or the source code, and it is provided in a form a
non-specialist can read alongside a plain-language key.

The chain either verifies or it does not, and the verification tool is provided
with the extract. That is a stronger commitment than a promise of accuracy,
because it can be checked by someone who does not trust the person making it.

\subsection{What happens if a rule in this section is broken}

Three breaches are the most likely, and each has a route that ends at a body
independent of the sponsor and is available to the participant without going
through the site.

\vspace{0.30em}

\noindent\begin{tabularx}{\textwidth}{L{4.5cm} L{4.0cm} Y}
\toprule
\textbf{Breach} & \textbf{First route} & \textbf{Independent route, always available to you} \\
\midrule
An adverse event affecting you is not reported to you & Site PI, on the 24-hour line & The IRB directly; the IRB may suspend approval \\
Your data is accessed by someone not on the read list & Site PI, who must log the request & The IRB directly, and the access is visible in the hash chain \\
A procedure proceeds under a software version you did not consent to & Sponsor-Investigator & The IRB and the Physical AI Safety Review Committee; the version registry is the evidence \\
\bottomrule
\end{tabularx}

\vspace{0.30em}

\noindent In each case the evidence needed to establish the breach is already in
the participant's own extract. That is the practical reason the audit trail is
worth having: it converts a dispute about what happened into a question about
what the record says.
